\documentclass[letterpaper,11pt,leqno]{article}
\usepackage{../../style/paper}
\bibliographystyle{../../style/paper}

% Enter document title to populate PDF metadata:
\hypersetup{pdftitle={Information Technology, Digital Media, and Digital Communications Committee}}

\begin{document}

% Enter title:
\title{Information Technology, Digital Media, and Digital Communications Committee}

% Enter author:
\author{Lacamas Watershed Council}

% Enter date:
\date{July, 2026}

% Enter document URL (can be commented out):
% \paperurl{https://lacamaswatershed.org}

\begin{titlepage}
\maketitle

% Enter abstract:

\end{titlepage}

% Enter main text:
\section{Executive Summary}\label{s:summary}

\paragraph{Purpose} The purpose of this document is to establish the operational framework for the LWC Information Technology, Digital Media, and Digital Communications Committee \footnote{Hereafter referred to as "The Committee"}.
This document will also establish the budget for the committee, describe chair and volunteer responsibilities, and designate flagship projects. 

IT, digital communications, and marketing have all been named under this committee as a matter of efficiency.
Though communications and marketing could be a separate committee, the majority of digital marketing tools will be provisioned by the IT officer and come from the IT budget. Therefore, treating LWC marketing, branding, and communications
under the IT umbrella reduces bureaucracy and opens a direct line of communication between the IT officer and LWC digital marketing volunteers.

\paragraph{Goals} The goal of The Committee is to support the LWC's vision \footnote{Establishing a safe, healthy and enjoyable watershed with a robust lake management plan safeguarding the future.} by providing efficient technological support, unified branding, 
and consistent, honest communication to the community and our partners. We will utilize technology and a modern digital presence to restore, protect, and advocate for our watersheds. We will never recommend a technology or media strategy that contradicts the LWC's mission and vision, and
we will never advocate for technology or media strategies that directly harm our local, regional, or global ecosystems.

\paragraph{Overview} The Committee chair must be a board member (ideally the Secretary or Information Officer). All other members can be LWC volunteers on or off The Board. The Committee has two branches: information technology and digital media and communications.
The information technology branch is responsible for managing our Google Workspace, databases, and other digital resources. The digital media and communications branch is responsible for creating lwc branding, creating our content calendar, taking and editing photos, and sending newsletters.
On an annual basis, The Committee will evaluate current technology and branding trends and propose adding or removing certain technologies to The Board for approval. Also on an annual basis, the IT branch will conduct an audit of our Google Workspace in order to better organize and de-clutter the Google Drive.

\paragraph{Budget} The total proposed annual budget is \$1,210 with \$384 allocated to the IT branch and \$826 allocated to the digital media and communications branch.

\section{Approval}\label{s:approval}
After deliberation the board will appoint a committee chair and vote to approve this document along with any amendments. Approving this document funds the committee through 2027.
This includes funding a prorated budget of \$504 for August-December 2026.

\section{Committee Structure}\label{s:structure}

\subsection{Information Technology Branch}\label{s:itBranch}
The information technology branch is responsible for managing the LWC's Google Workspace, databases, and other digital resources. The IT branch purchases, sets up, and maintains all LWC digital tools so that 
the organization can succeed in a modern digital environment.

\paragraph{Budget} The proposed budget is \$384. The IT branch needs a total of \$288 minimum to meet our website and domain name needs. For an additional \$96 we 
should add a password manager for 2 individuals (President and IT Chair). All other IT services are either free or covered under our AWS nonprofit credits.
\paragraph{Tools and cost}
\begin{itemize}
    \item Google Workspace: Free at our current usage
    \item Amazon Web Services (database, data dashboard): \$1,000 in nonprofit credits.
    \item Domain name: \$11.64
    \item Squarespace (website): \$276\
    \item Password manager (proposed): \$96
\end{itemize}

\subsection{Digital Media and Communications Branch}\label{s:dmBranch}
The digital media and communications branch is responsible for creating LWC branding, managing the content calendar, taking and editing photos, and sending newsletters.
The digital media branch also covers printed media including stickers, fliers, t-shirts, door hangers, and posters. However, it is the recommendation of the committee to use
printed media sparingly. This is a cost saving measure and also reinforces our stance as an environmentally sustainable nonprofit.

\paragraph{Budget} The proposed budget is \$826. The DM branch needs a total of \$156 annually to cover mailchimp. In addition, we propose an additional \$150 for social media advertising 
and \$420 for Adobe Creative Cloud Pro to support photo editing and graphic design.
\paragraph{Tools and cost}
\begin{itemize}
    \item Mailchimp: \$156
    \item Social media advertising (proposed): \$150
    \item Adobe Creative Cloud Pro (proposed): \$420
    \item Printed media (proposed): \$100 (primarily for up front t-shirt costs that will be recouped)
\end{itemize}

\section{Roles and Responsibilities}\label{s:roles}

\subsection{Committee Chair}\label{s:chair}

The committee chair must be a board member and, ideally, is the LWC Secretary or Information Officer. The chair is responsible for maintaining the committee's budget, organizing volunteers, and ensuring that the committee's goals are met. The chair will also serve as the primary point of contact between the committee and The Board. However,
most importantly, the chair should have domain knowledge of information technology and information security, given that the committee is responsible for managing the LWC's private digital resources. Other duties of the chair include, but are not limited to:
\begin{itemize}
\item Orchestrating regular communications with IT and digital media volunteers with the latest board decisions.
\item Managing the LWC's Google Workspace, including user accounts, permissions, and security settings.
\item Managing the LWC's water quality database and establishing data sharing policies with other organizations.
\item Ensuring consistent uptime of the LWC's website.
\item Finding and implementing the correct IT tools to support the LWC's mission and vision.
\item Finding and implementing the correct digital media tools to ensure consistent branding and messaging across all LWC digital platforms.
\end{itemize}

\subsection{Other IT \& Digital Media Volunteers}\label{s:dmVolunteer}

All digital media volunteers are responsible for creating content for the LWC's digital platforms, including social media, newsletters, and the website. When appropriate, they will
also create printed media. Some common volunteer positions that fall under this category include:

\paragraph{Communications \& Outreach Volunteer} Responsible for creating the LWC's content calendar, brand guide, and social media strategy. This volunteer will also be responsible for creating and sending newsletters, and managing the LWC's newsletter contacts database.
The C\&O volunteer is the board's voice to the public and will provide guidance to the board on how to best communicate LWC operations to the widest possible audience. 

\paragraph{Photographer \& Videographer Volunteer} Responsible for taking and editing photos and videos for the LWC's digital communications and social media. This volunteer should understand how to use
the basic functions of a digital camera and have a basic understanding of photo and video editing software.

\paragraph{Social Media Volunteer} Responsible for creating and posting content to the LWC's social media platforms and adhering to the content calendar. This volunteer should have a basic understanding of social media platforms and how to drive social engagement.

\paragraph{Website Volunteer} Responsible for maintaining the LWC's website and ensuring that it is up to date with the latest LWC information. This volunteer should have a basic understanding of web development and content management systems (like squarespace or WordPress).
% To include a figure, place the image file in figures/ and use:
%
% \begin{figure}[t]
% \includegraphics[width=\textwidth]{figures/filename}
% \caption{Caption text.}
% \note{Note text.}
% \label{f:label}\end{figure}

\section{Amendments and Adding Projects}
\subsection{Amending this Document}

\paragraph{Major Amendments} A major amendment is a significant change to the document that affects its overall structure or purpose. This includes changing goals, adding new projects, and adding or modifying defined roles.
To make a major amendment, a board member must submit a proposed amendment for review by the board. The proposed amendment must be approved according to the LWC's bylaws (including quorum requirements per Article 3.6).
Once approved, the proposed amendment will be added to this document within thirty (30) days by the committee chair, and the document will be re-published and distributed to the relevant stakeholders.
For example, if an amendment changes the field operations of a particular project, the project manager and volunteers of that project will be notified of the change.
The committee chair is responsible for maintaining this document and ensuring that it is up to date with the latest amendments.

\paragraph{Minor Amendments} A minor amendment is a small change to the document that does not significantly affect its overall structure or purpose. This includes correcting typos, updating project statuses, and other small changes.
The committee chair makes the initial determination of whether an amendment is major or minor. To make a minor amendment, the committee chair will initiate the amendment(s) and notify all board members of the changes at a board meeting or by email to the full board. The amendment will be incorporated into this document within thirty (30) days of notification. At any point, a board member can request that a minor amendment be reviewed as a major amendment, at which point the amendment will be held pending the outcome of that review.

\bibliography{refs/paper}
\end{document}
